\lecture{11}{Tue 27 Sep 2022 12:04}{}

\begin{theorem}[Cauchy Theorem]
	If $f$ is analytic in $\Omega$ and $\mathcal{D}$ is a region, $\partial \mathcal{D}$ is piecewise smooth, $\mathcal{D} \cup \partial \mathcal{D} \subset \Omega$, Then \[
		\int _{\partial \mathcal{D}}f(z) dz = 0
	.\] 
\end{theorem}

\begin{example}
	if $f$ is rational and has no poles in $\Omega$, consider the path integral \[
		\int _{\partial \mathcal{D}} \frac{f(z)}{z - a}dz
	.\] 
	Where the partial fraction decomposition is (assume $f(a) \neq  0),$
	\[
		\frac{f(z)}{z - a} = \frac{f(a)}{z - a}+ g
	\] 
	where $g$ is analytic in $\Omega$.

	Now $\int_{\partial \mathcal{D}} g(z) dz= 0$ by Cauchy's Theorem. Therefore: \[
		\int_{\partial \mathcal{D}}\frac{f(z)}{z - a}dz 
		= \int_{\partial \mathcal{D}}\frac{f(a)}{z - a}dz
		= 2 \pi i f(a)
	.\] 
\end{example}

The next theorem generalizes the example to all analytic functions.

\begin{corollary}[Cauchy Integral Formula]
	If $f$ is analytic in $\Omega$ and $\mathcal{D}$ is a region, $\partial \mathcal{D}$ is piecewise smooth, $\mathcal{D} \cup \partial \mathcal{D} \subset \Omega$, Then for all $a$ in the interior of $\mathcal{D}$,
	\[
		f(a) = \frac{1}{2 \pi i } \int_{\partial \mathcal{D}} \frac{f(z)}{z-a}dz
	\]
\end{corollary}
\begin{proof}
	We remove some disk in $\mathcal{D}$ centered at $a$, say $\mathcal{D}_{\varepsilon}$.
	\[
		\mathcal{D}_\varepsilon = \mathcal{D} \setminus \{z: |z-a| \le \varepsilon\} 
	\] 
	\[
		\partial \mathcal{D}_\varepsilon = \partial \mathcal{D} - C_\varepsilon
	\] 
	Where $C_\varepsilon$ is oriented around the small disk counterclockwise. Now we have removed the singularity $a$, hence
	$\frac{f(z)}{z-a}$ is analytic in $\mathcal{D}_\varepsilon$. By Cauchy Theorem, \[
		\int _{\partial \mathcal{D}_\varepsilon} \frac{f(z)}{z-a}dz = 0
	.\] 
	Now \[
		\int _{\partial\mathcal{D}} - \int _{C_\varepsilon} = \int _{\partial\mathcal{D}_\varepsilon} = 0
	.\] 
	Therefore
	\[
		\int _{C_\varepsilon}\frac{f(z)}{z-a}dz
		=\int _{\partial \mathcal{D}}\frac{f(z)}{z-a}dz 
	.\] 
	By the integral $\int _{C_\varepsilon}\frac{1}{z-a}dz = 2 \pi i$,
	\[
		f(a) = \frac{1}{2 \pi i} \int _{C_\varepsilon} \frac{f(a)}{z - a}dz
	.\] 
Now estimate the difference:
\begin{align*}
	&\left| \frac{1}{2 \pi i}\int _{C_\varepsilon}\frac{f(a)}{z-a}dz 
	- \frac{1}{2 \pi i} \int _{C_\varepsilon} \frac{f(z)}{z-a} dz\right| \\
	&\le 2 \pi \varepsilon \frac{1}{2\pi} \frac{\max |f(z) - f(a)|}{\varepsilon}\\
	&\to 0 
.\end{align*}
Hence \[
	f(a) = \frac{1}{2 \pi i} \int_{\partial \mathcal{D}}\frac{f(z)}{z-a}
.\]
\end{proof}

\begin{remark}
	Boundary determines interior points. This is similar to Dirichlet Problem and harmonic functions.
\end{remark}

\begin{lemma}
	Let $g$ be continuous on the contour $\Gamma$, and for each $z$ not on $\Gamma$ set
	\[
		G(z) = \int_\Gamma \frac{g(\xi)}{\left( \xi - z \right) ^n} d \xi
	.\] 
	Then $G$ is analytic at each point not on $\Gamma$, and its derivative is
	\[
		G'(z) = n \int _\Gamma \frac{g(\xi)}{(\xi - z)^{n+1}} d \xi
	.\] 
\end{lemma}
Note: we do not assume $g$ is analytic, or $\Gamma$ is closed.
\begin{proof}
	\begin{align*}
		G'(z)
		&= \lim_{h \to 0} \frac{G(z + h) - G(z)}{h} \\
		&= \lim_{h \to 0} \frac{1}{h} \int _\gamma \frac{g(\xi)}{(\xi - z - h)^n} - \frac{g(\xi)}{(\xi - z)^n} d \xi \\
		&= \lim_{h \to 0} \int _\gamma \frac{g(\xi) \left( (\xi - z)^n - (\xi - z - h)^n \right) }{(\xi - z - h)^n (\xi - z)^n} \\
		&= \lim_{h \to 0} \int _\gamma \frac{g(\xi) \left( (\xi - z) - (\xi - z - h) \right)\left( (\xi - z)^{n - 1} + \ldots + (\xi - z- h)^{n - 1} \right)  }{(\xi - z - h)^n (\xi - z)^n} \\
		&= \int _\gamma \frac{g(\xi) n (\xi - z)^{n - 1}}{(\xi - z)^{2 n}} \\
		&= n \int _\gamma \frac{g(\xi)}{(\xi - z)^{n + 1}}
	.\end{align*}
\end{proof}


\begin{corollary}[Cauchy's Differentiation Formula]
	\textbf{Derivatives are analytic}, so differentiable functions are infinitely differentiable. Moreover,
\[
	\frac{d^n}{dz^n} f(z) = \frac{n!}{2 \pi i} \int _{\partial \mathcal{D}}\frac{f(\zeta)}{(\zeta - z)^{n+1}} d\zeta
\]
for all $z$ in the interior of $\mathcal{D}$.
\end{corollary}  
\begin{proof}
	Induct on $n$. The base case $n = 0$ is given by \logseq{Cauchy Integral Formula}. The induction step is given by the previous lemma.
\end{proof}

\begin{theorem}[Morera's Theorem]
	If $g$ is continuous in a domain $D$ and if for any closed contour $\Gamma$ we have
	\[
		\int _\Gamma f(z) dz = 0
	.\] 
	Then $f$ is analytic in $D$.
\end{theorem}
\begin{proof}
	Integral over each closed path is $0$, so $f$ has a primitive. Derivaties are analytic, so $f$ is analytic.
\end{proof}


\begin{property}[Average Property of Analytic Functions]
	
Let $\mathcal{D} = \{z: |z - a| > r\}$. Then $f(a) = \frac{1}{2 \pi i}\int _\gamma \frac{f(z)}{z - a} dz$.

We parameterize $\gamma$. Let $\gamma(t) = a + r e^{i t}, 0\le t\le 2 \pi$. $\gamma'(t) = i r e^{it}$. Now
\begin{align*}
	f(a) &= \frac{1}{2 \pi i}\int _0^{2 \pi}\frac{f(a + r e ^{it })}{r e^{it }} i r e^{it } dt\\
	     &= \frac{1}{2 \pi} \int _0^{2 \pi} f(a + r e^{it }) dt
.\end{align*}
This is the average of $f$ along the curve.

By splitting real and imaginary part, we conclude that this property holds for real and imaginary parts separately. Hence \logseq{Harmonic function}s have average property.
\end{property}

Suppose $g: \mathcal{D} \to \mathbb{R}$ is harmonic. We say $a$ is a point of (local, non-strict) maximum if $g(a) \ge g(z)$ for $z$ in some neighbourhood  of $a$.

\begin{theorem}[Maximum Principle]
	\begin{enumerate}
		\item If a function $u$ harmonic in $\mathcal{D}$ has a (local, nonstrict) maximum at some point $A \in \mathcal{D}$, then $u$ is constant.
		\item Minimum Principle: same for mimima of harmonic functions.
	\end{enumerate}
\end{theorem}

\begin{corollary}
Solution to \logseq{Dirichlet Problem } for bounded region must be unique.
\end{corollary}
\begin{proof}
	Suppose $u_1, u_2$ are two solutions, then $u_1-u_2$ is harmonic, and have boundary value zero. 

	If $u_1-u_2 \neq 0$, there must be a maximum or minimum in domain. ($u_1-u_2$ is continuous on a close and bounded region, and continuous functions must achieve extremum on compact domain). Then by maximum principle, $u_1 = u_2$.
\end{proof}

\begin{theorem}[Maximum Modulus Principle]
	If a function $f$ is analytic in $\mathcal{D}$, and $|f(z)|$ has a (local, nonstrict) maximum in $\mathcal{D}$ then $f$ is constant.
\end{theorem}
\begin{proof}
	We have $f(a) = \frac{1}{2 \pi}\int _0^{2 \pi} f(a + r e^{i t})d t$. Suppose $|f(z)|$ has a maximum in $\mathcal{D}$, we take $r  $ small enough so that $|f(z+re^{it })| \le f(a)$.

	On the other hand, \[
		|f(a)| \le  \frac{1}{2 \pi} \int _0^{2 \pi} |f(a + r e^{i t})|d t
	.\] 
	These imply that $|f(a + r e^{i t})|$ is constant.
\end{proof}
We do not have ``minimum modulus principle'' in the same form. To obtain it we need to replace $f$ by $\frac{1}{f}$.

\begin{example}
	Consider $\sin(z)$. The real version has infinitely many extrema. However the complex version is harmonic, and the previous maxima become saddle points. 

	\logseq{Harmonic functions} cannot have maxima if non-constant. They can only have saddle points.
\end{example}

\begin{corollary}
	If $\mathcal{D}$ bounded, $\overline{\mathcal{D}} $ is compact, $f$ analytic in $\mathcal{D}$, continuous in $\overline{\mathcal{D}} $. Then $|f|$ has a global maximum in $\mathcal{D}$. Furthermore, this global maximum must be in $\partial \mathcal{D}$.
\end{corollary}

\begin{corollary}
	Under the same assumptions in the last corollary, if $f$ bounded by $M$ in $\partial \mathcal{D}$, then $f$ is bounded by $M$ in $\mathcal{D}$.
\end{corollary}
